
\par\noindent\textit{引理~\ref{lem:7-2-310} 的证明。}
令
\[
\begin{gathered}
C'_0=C([0,1]_s\times L'_0,\lambda'_0),\qquad
C'=C([0,1]_s\times L'_0,(\lambda_0,\lambda'_0)),\\
\widehat{\mathfrak C}=\widehat{\mathfrak C}_a,\qquad
\widehat{\mathfrak C}_{s=s_0}=\mathfrak C_{a,r},\qquad
\widehat{\mathfrak C}_{r=r_0}
=C([0,1]_s\times L'_1,\lambda'_1).
\end{gathered}
\]
最后一个端面模型按 (7.2.307) 取定。\translatornote{原书此处将
$\widehat{\mathfrak C}_{r=r_0}$
写成
$C_a([0,1]_s\times L'_1,(\lambda'_0,\lambda'_1))$；
但 (7.2.307) 已将其定义为
$C([0,1]_s\times L'_1,\lambda'_1)$，
故依前式定型。}
定义
\[
\begin{aligned}
\mathfrak H^{r=0}
&=j_{[0,1]_s}\circ
\mathfrak f^{\operatorname{id}\times\phi'}\circ\mathfrak e_a,\\
\mathfrak H^{r=1}
&=j_{[0,1]_s}\circ\mathfrak E_a\circ
\mathfrak f^{\operatorname{id}\times\phi'}_{\mathrm{ext}},\\
\mathfrak H^{s=s_0}
&=\mathcal F_a^{s=s_0}\circ\operatorname{Eval}_{s=s_0}.
\end{aligned}
\]
它们分别是 (7.2.311.1)--(7.2.311.3) 的右端。
另按 (7.2.311.4) 定义
\[
\mathfrak H_0
=\operatorname{Incl}_r\circ j_{[0,1]_s}
\circ\mathfrak f^{\operatorname{id}\times\phi'}.
\]
最后，取引理~\ref{lem:7-2-300}
的 (7.2.301.5) 中的包含映射
\[
\mathfrak h\colon
C([0,1]_s\times L'_0,\lambda'_0)
\longrightarrow
C([0,1]_s\times L'_0,(\lambda_0,\lambda'_0)).
\]

我们验证 (7.2.313.1)--(7.2.313.3)。
结合 (7.2.305.4)、(7.2.305.2)
和 (7.2.301.1)，得到
\[
\begin{aligned}
&\operatorname{Eval}_{s=s_0}\circ j_{[0,1]_s}
\circ\mathfrak f^{\operatorname{id}\times\phi'}
\circ\mathfrak e_a\\
&\qquad
=j\circ\mathfrak f^{\phi'}
\circ\operatorname{Eval}_{s=s_0}\circ\mathfrak e_a\\
&\qquad
=j\circ\mathfrak f^{\phi'}\circ j_{s=s_0}
\circ\operatorname{Eval}_{s=s_0}.
\end{aligned}
\]
另一方面，
将 $\operatorname{Eval}_{s=s_0}$
与 (7.2.306.1) 复合，可得
\[
\operatorname{Eval}_{r=0}\circ\mathcal F_a^{s=s_0}
\circ\operatorname{Eval}_{s=s_0}
=j\circ\mathfrak f^{\phi'}\circ j_{s=s_0}
\circ\operatorname{Eval}_{s=s_0}.
\]
故
\begin{equation}
\operatorname{Eval}_{s=s_0}\circ\mathfrak H^{r=0}
=\operatorname{Eval}_{r=0}\circ\mathfrak H^{s=s_0}.
\tag{7.2.323}\label{eq:7-2-323}
\end{equation}
这一次对 $r=1$ 应用
(7.2.305.4)、(7.2.305.3)、
(7.2.301.2) 和 (7.2.306.2)，
同理证明
\begin{equation}
\operatorname{Eval}_{s=s_0}\circ\mathfrak H^{r=1}
=\operatorname{Eval}_{r=1}\circ\mathfrak H^{s=s_0}.
\tag{7.2.324}\label{eq:7-2-324}
\end{equation}
结合 (7.2.323) 和 (7.2.324)，
即证明了 (7.2.313.1)。

为验证 (7.2.313.2)，计算得
\begin{equation}
\operatorname{Eval}_{r=r_0}\circ\operatorname{Incl}_r
\circ j_{[0,1]_s}\circ
\mathfrak f^{\operatorname{id}\times\phi'}
=j_{[0,1]_s}\circ
\mathfrak f^{\operatorname{id}\times\phi'}.
\tag{7.2.325}\label{eq:7-2-325}
\end{equation}
再由 (7.2.305.1)，在
$C([0,1]_s\times L'_0,\lambda'_0)$
上有\translatornote{原书此处的定义域写作
$C_a([0,1]_s\times L'_0,\lambda'_0)$；
证明开头的 $C'_0$ 和 (7.2.305.1)
使用的均为无下标的 $C$。}
\[
j_{[0,1]_s}\circ
\mathfrak f^{\operatorname{id}\times\phi'}
\circ\mathfrak e_a
=j_{[0,1]_s}\circ
\mathfrak f^{\operatorname{id}\times\phi'}.
\]
这与 (7.2.325) 一致。
再利用 (7.2.301.6) 和 (7.2.305.1)，
可知在
$C([0,1]_s\times L'_0,\lambda'_0)$
上还有
\[
j_{[0,1]_s}\circ\mathfrak E_a\circ
\mathfrak f^{\operatorname{id}\times\phi'}_{\mathrm{ext}}
=j_{[0,1]_s}\circ
\mathfrak f^{\operatorname{id}\times\phi'}.
\]
这也与 (7.2.325) 一致，
故 (7.2.313.2) 得证。

最后验证 (7.2.313.3)。
由 (7.2.308)、(7.2.305.4)
及 (7.2.301.1)，计算得
\begin{equation}
\begin{aligned}
&\operatorname{Eval}_{s=s_0}\circ
\operatorname{Incl}_r\circ j_{[0,1]_s}
\circ\mathfrak f^{\operatorname{id}\times\phi'}\\
&\qquad
=\operatorname{Incl}_r\circ j\circ\mathfrak f^{\phi'}
\circ\operatorname{Eval}_{s=s_0}.
\end{aligned}
\tag{7.2.326}\label{eq:7-2-326}
\end{equation}
另一方面，
由于
\[
\operatorname{Eval}_{s=s_0}
\bigl(C([0,1]_s\times L'_0,\lambda'_0)\bigr)
\subset C(L'_0,\lambda'_0),
\]
(7.2.306.3) 蕴涵
在 $C([0,1]_s\times L'_0,\lambda'_0)$ 上有
\[
\mathcal F_a^{s=s_0}\circ\operatorname{Eval}_{s=s_0}
=\operatorname{Incl}_r\circ j\circ\mathfrak f^{\phi'}
\circ\operatorname{Eval}_{s=s_0}.
\]
这与 (7.2.326) 一致，
从而证明了 (7.2.313.3)。
现在可应用命题~\ref{prop:7-2-312}，
引理~\ref{lem:7-2-310} 随之成立。\(\square\)
